\documentclass[12pt, a4paper]{article}

% Essential Math Packages
\usepackage{amsmath}
\usepackage{amssymb}
\usepackage{amsthm}
\usepackage{geometry}

% Page Setup and Margins
\geometry{a4paper, margin=1in}

% Header and Footer Configuration
\usepackage{fancyhdr}
\pagestyle{fancy}
\fancyhf{}
\lhead{From Calculus to Cohomology}
\chead{Homework Set \# 2}
\rhead{Name: Zijian Zhu}
\lfoot{Instructor: Yuan Gao}
\rfoot{Nanjing University}
\cfoot{\thepage}

% Theorem and Problem Environments
\theoremstyle{definition}
\newtheorem{problem}{Problem}

% Custom Solution Environment (uses a black square for the QED symbol)
\newenvironment{solution}
  {\renewcommand\qedsymbol{$\blacksquare$}\begin{proof}[Solution]}
  {\end{proof}}

% Custom math operators
\DeclareMathOperator{\curl}{curl}
\DeclareMathOperator{\divg}{div}
\DeclareMathOperator{\grad}{grad}
\DeclareMathOperator{\im}{im}
\DeclareMathOperator{\sgn}{sgn}

\begin{document}

% Academic Integrity / Collaboration Declaration
\noindent \textbf{Collaborators:} \textit{None}
\vspace{0.8cm}

% --- Begin Homework Problems ---

\begin{problem}
Let $\omega$ be the 1-form on $\mathbb{R}^3$ given by 
$$\omega = x^2 y \, dx + \sin(xz) \, dy + e^{yz} \, dz$$
Compute its exterior derivative $d\omega$, grouping your final answer in terms of the standard basis 2-forms ($dy \wedge dz$, $dx \wedge dz$, and $dx \wedge dy$).
\end{problem}

\begin{solution}
\[
d\omega = d(x^{2}y) \wedge dx + d(\sin(xz)) \wedge dy + d(e^{yz}) \wedge dz.
\]
We compute each term:
\[
d(x^{2}y) = 2xy\,dx + x^{2}\,dy \quad\Rightarrow\quad (2xy\,dx + x^{2}\,dy) \wedge dx = -x^{2}\,dx\wedge dy.
\]
\[
d(\sin(xz)) = z\cos(xz)\,dx + x\cos(xz)\,dz \quad\Rightarrow\quad (z\cos(xz)\,dx + x\cos(xz)\,dz) \wedge dy = z\cos(xz)\,dx\wedge dy - x\cos(xz)\,dy\wedge dz.
\]
\[
d(e^{yz}) = z e^{yz}\,dy + y e^{yz}\,dz \quad\Rightarrow\quad (z e^{yz}\,dy + y e^{yz}\,dz) \wedge dz = z e^{yz}\,dy\wedge dz.
\]
Summing, we obtain
\[
d\omega = \bigl(-x^{2} + z\cos(xz)\bigr)\,dx\wedge dy + \bigl(z e^{yz} - x\cos(xz)\bigr)\,dy\wedge dz.
\]
\end{solution}

\vfill

\begin{problem}
Let $U = \mathbb{R}^2$ with coordinates $(u,v)$ and $V = \mathbb{R}^3$ with coordinates $(x,y,z)$. Consider the smooth map $F: U \to V$ given by $F(u,v) = (u^2, v^2, uv)$. Let $\omega$ be the 2-form on $V$ defined by:
$$\omega = x \, dy \wedge dz + z \, dx \wedge dy$$
Compute the pullback $F^*\omega$, expressing your final answer entirely in terms of $u, v$ and $du \wedge dv$.
\end{problem}

\begin{solution}
\[
F^{*}\omega = (x\circ F)\,d(y\circ F)\wedge d(z\circ F) + (z\circ F)\,d(x\circ F)\wedge d(y\circ F).
\]
Here \(x=u^{2}\), \(y=v^{2}\), \(z=uv\). Then
\[
d(y\circ F)=2v\,dv,\qquad d(z\circ F)=u\,dv+v\,du,\qquad d(x\circ F)=2u\,du.
\]
First term:
\[
u^{2}\cdot (2v\,dv)\wedge (u\,dv+v\,du)=2u^{2}v\,(u\,dv\wedge dv + v\,dv\wedge du)=2u^{2}v^{2}\,dv\wedge du = -2u^{2}v^{2}\,du\wedge dv.
\]
Second term:
\[
uv\cdot (2u\,du)\wedge (2v\,dv)=uv\cdot 4uv\,du\wedge dv = 4u^{2}v^{2}\,du\wedge dv.
\]
Adding gives
\[
F^{*}\omega = (-2u^{2}v^{2}+4u^{2}v^{2})\,du\wedge dv = 2u^{2}v^{2}\,du\wedge dv.
\]
\end{solution}

\vfill\newpage

\begin{problem}
This problem verifies the Graded Leibniz Rule for specific forms. Let $\alpha = x \, dy + y \, dz$ and $\beta = z \, dx + x \, dy$ be 1-forms on $\mathbb{R}^3$.
\begin{enumerate}
    \item[(a)] Compute $d\alpha$ and $d\beta$.
    \item[(b)] Compute the wedge product $\alpha \wedge \beta$ and then compute its exterior derivative $d(\alpha \wedge \beta)$.
    \item[(c)] Because $\alpha$ is a 1-form ($k=1$), the Graded Leibniz Rule states that $d(\alpha \wedge \beta) = d\alpha \wedge \beta - \alpha \wedge d\beta$. Compute $d\alpha \wedge \beta - \alpha \wedge d\beta$ and verify that it equals your result from part (b).
\end{enumerate}
\end{problem}

\begin{solution}

(a)\[
d\alpha = dx\wedge dy + dy\wedge dz,\qquad d\beta = dz\wedge dx + dx\wedge dy.
\]

(b)\[
\alpha\wedge\beta = (x\,dy + y\,dz)\wedge(z\,dx + x\,dy) = -xz\,dx\wedge dy - yz\,dx\wedge dz - xy\,dy\wedge dz.
\]
Now compute the exterior derivative:
\[
d(-xz\,dx\wedge dy) = -(z\,dx + x\,dz)\wedge dx\wedge dy = -x\,dz\wedge dx\wedge dy = -x\,dx\wedge dy\wedge dz,
\]
\[
d(-yz\,dx\wedge dz) = -(z\,dy + y\,dz)\wedge dx\wedge dz = -z\,dy\wedge dx\wedge dz = z\,dx\wedge dy\wedge dz,
\]
\[
d(-xy\,dy\wedge dz) = -(y\,dx + x\,dy)\wedge dy\wedge dz = -y\,dx\wedge dy\wedge dz.
\]
Summing: \(d(\alpha\wedge\beta) = (-x+z-y)\,dx\wedge dy\wedge dz = (z-x-y)\,dx\wedge dy\wedge dz.\)

(c)\[
d\alpha\wedge\beta = (dx\wedge dy + dy\wedge dz)\wedge(z\,dx + x\,dy) = z\,dy\wedge dz\wedge dx = z\,dx\wedge dy\wedge dz,
\]
\[
\alpha\wedge d\beta = (x\,dy + y\,dz)\wedge(dz\wedge dx + dx\wedge dy) = x\,dy\wedge dz\wedge dx + y\,dz\wedge dx\wedge dy = x\,dx\wedge dy\wedge dz + y\,dx\wedge dy\wedge dz = (x+y)\,dx\wedge dy\wedge dz.
\]
Thus \(d\alpha\wedge\beta - \alpha\wedge d\beta = (z-x-y)\,dx\wedge dy\wedge dz\), matching part (b).
\end{solution}

\vfill\newpage

\begin{problem}
Let $f$ and $g$ be smooth scalar functions ($0$-forms) on an open set $U \subseteq \mathbb{R}^n$.
\begin{enumerate}
    \item[(a)] Using only the intrinsic axioms of the exterior derivative (the Graded Leibniz Rule and the square-zero property $d^2 = 0$), mathematically prove that $d(f \, dg) = df \wedge dg$.
    \item[(b)] Using the same axioms, what is $d(df \wedge dg)$? Briefly justify your answer.
\end{enumerate}
\end{problem}

\begin{solution}

(a)By the graded Leibniz rule with \(\alpha=f\) (0‑form) and \(\beta=dg\) (1‑form):
\[
d(f\wedge dg) = df\wedge dg + (-1)^{0}f\wedge d(dg) = df\wedge dg + f\wedge d^{2}g = df\wedge dg,
\]
since \(d^{2}g=0\).

(b)With \(\alpha=df\) (1‑form) and \(\beta=dg\) (1‑form):
\[
d(df\wedge dg) = d^{2}f\wedge dg - df\wedge d^{2}g = 0.
\]
\end{solution}

\vfill\newpage

\begin{problem}
One of the most important properties of the pullback is that it commutes with the exterior derivative: $d(F^*\omega) = F^*(d\omega)$. We will verify this using the map $F$ and the 2-form $\omega$ from Problem 2.
\begin{enumerate}
    \item[(a)] First, calculate the exterior derivative $d\omega$ on $V = \mathbb{R}^3$. What must the pullback $F^*(d\omega)$ be purely based on dimensional constraints on $U = \mathbb{R}^2$?
    \item[(b)] Take your final pulled-back form $F^*\omega$ from Problem 2 and compute its exterior derivative $d(F^*\omega)$. Does it match your answer from part (a)?
\end{enumerate}
\end{problem}

\begin{solution}

(a)First compute \(d\omega\) on \(V=\mathbb{R}^{3}\).
\[
\omega = x\,dy\wedge dz + z\,dx\wedge dy \quad\Rightarrow\quad d\omega = dx\wedge dy\wedge dz + dz\wedge dx\wedge dy = 2\,dx\wedge dy\wedge dz.
\]
Since \(U=\mathbb{R}^{2}\) is two‑dimensional, any 3‑form pulls back to zero, so \(F^{*}(d\omega)=0\).

(b)From Problem 2, \(F^{*}\omega = 2u^{2}v^{2}\,du\wedge dv\). Then
\[
d(F^{*}\omega) = d(2u^{2}v^{2})\wedge du\wedge dv = (4u v^{2}\,du + 4u^{2}v\,dv)\wedge du\wedge dv = 0,
\]
which matches \(F^{*}(d\omega)\).
\end{solution}

\vfill\newpage

\begin{problem}
Let $V = \mathbb{R}^2 \setminus \{(0,0)\}$ and consider the angle form $\omega = \frac{-y}{x^2+y^2}dx + \frac{x}{x^2+y^2}dy$. Let $U = (0, 2\pi)$ with coordinate $t$, and let $\iota: U \to V$ be the parameterization of the unit circle given by $\iota(t) = (\cos t, \sin t)$.
\begin{enumerate}
    \item[(a)] Compute the pullback $\iota^*\omega$, expressing your answer entirely in terms of $t$ and $dt$.
    \item[(b)] Use your answer to evaluate the integral $\int_0^{2\pi} \iota^*\omega$.
\end{enumerate}
\end{problem}

\begin{solution}

(a)Compute \(\iota^{*}\omega\).
Substituting \(x=\cos t,\;y=\sin t\):
\[
\iota^{*}\omega = \frac{-\sin t}{1}\,d(\cos t) + \frac{\cos t}{1}\,d(\sin t) = -\sin t(-\sin t\,dt) + \cos t(\cos t\,dt) = (\sin^{2}t+\cos^{2}t)\,dt = dt.
\]
\end{solution}

(b)Evaluate \(\int_{0}^{2\pi}\iota^{*}\omega\).
\[
\int_{0}^{2\pi} dt = 2\pi.
\]

\vfill\newpage

\begin{problem}
Let $U \subseteq \mathbb{R}^3$ and consider an arbitrary 3-form $\omega = A(x,y,z) \, dx \wedge dy \wedge dz$.
\begin{enumerate}
    \item[(a)] Compute the exterior derivative $d\omega$. 
    \item[(b)] Based on your result and the algebraic properties of the exterior algebra $\Lambda^k(\mathbb{R}^n)^*$, briefly explain why \emph{every} $n$-form on an $n$-dimensional open set must be closed.
\end{enumerate}
\end{problem}

\begin{solution}

(a)\[
d\omega = dA\wedge dx\wedge dy\wedge dz.
\]
Since \(dA\) is a 1‑form, the wedge product is a 4‑form, which must be zero on a 3-dimensional space. Hence \(d\omega=0\).

(b)Explain why every \(n\)-form on an \(n\)-dimensional open set is closed.
The exterior derivative raises the degree by one, but there are no non‑zero \((n+1)\)-forms in an \(n\)-dimensional manifold. Therefore \(d\omega=0\) for any \(n\)-form \(\omega\).
\end{solution}

\vfill\newpage

\begin{problem}
Let $U = \mathbb{R}^3 \setminus \{(0,0,0)\}$. Consider the 1-form:
$$\omega = \frac{x \, dx + y \, dy + z \, dz}{(x^2+y^2+z^2)^{3/2}}$$
\begin{enumerate}
    \item[(a)] Compute $d\omega$ to show that $\omega$ is a closed form.
    \item[(b)] Show that $\omega$ is also an exact form by explicitly verifying that $\omega = df$ for the function $f(x,y,z) = -(x^2+y^2+z^2)^{-1/2}$.
    \item[(c)] In Lecture 1, we learned that the ``angle form'' on the punctured plane $\mathbb{R}^2 \setminus \{(0,0)\}$ is closed but \emph{not} exact. Why does the exactness of $\omega$ on the punctured space $\mathbb{R}^3 \setminus \{(0,0,0)\}$ not contradict this topological principle? Briefly reference the relevant cohomology group from Lecture 1.
\end{enumerate}
\end{problem}

\begin{solution}

(a)Note that \(\omega = df\) with \(f=-(x^{2}+y^{2}+z^{2})^{-1/2}\) (see part (b)). Hence \(d\omega = d^{2}f = 0\).

(b)Let \(r = \sqrt{x^{2}+y^{2}+z^{2}}\). Then \(f = -r^{-1}\) and
\[
df = d(-r^{-1}) = \frac{1}{r^{2}}\,dr = \frac{1}{r^{2}}\cdot\frac{x\,dx+y\,dy+z\,dz}{r} = \frac{x\,dx+y\,dy+z\,dz}{r^{3}} = \omega.
\]

(c)Why does exactness on \(\mathbb{R}^{3}\setminus\{0\}\) not contradict the non‑exactness of the angle form on \(\mathbb{R}^{2}\setminus\{0\}\)?
The first de Rham cohomology \(H^{1}(\mathbb{R}^{3}\setminus\{0\})\) is trivial because \(\mathbb{R}^{3}\setminus\{0\}\) is simply connected, whereas \(H^{1}(\mathbb{R}^{2}\setminus\{0\})\cong\mathbb{R}\) (generated by the angle form). Thus closed 1‑forms on \(\mathbb{R}^{3}\setminus\{0\}\) are always exact.
\end{solution}

\vfill\newpage

\begin{problem}
Consider the 1-form $\omega = (2xy) \, dx + (x^2 + 3z) \, dy + (3y) \, dz$ on $\mathbb{R}^3$.
\begin{enumerate}
    \item[(a)] Compute $d\omega$ to verify that $\omega$ is a closed form.
    \item[(b)] Since $\mathbb{R}^3$ is star-shaped with respect to the origin, the Poincar\'e Lemma guarantees $\omega$ is exact. Apply the explicit homotopy operator formula $h\omega$ (for $k=1$) to find a function $f = h\omega$ such that $df = \omega$.
    \item[(c)] Verify your answer by computing the total differential $df$.
\end{enumerate}
\end{problem}

\begin{solution}

(a)
\[
d\omega = d(2xy)\wedge dx + d(x^{2}+3z)\wedge dy + d(3y)\wedge dz.
\]
\[
d(2xy)=2y\,dx+2x\,dy \;\Rightarrow\; (2y\,dx+2x\,dy)\wedge dx = -2x\,dx\wedge dy.
\]
\[
d(x^{2}+3z)=2x\,dx+3\,dz \;\Rightarrow\; (2x\,dx+3\,dz)\wedge dy = 2x\,dx\wedge dy -3\,dy\wedge dz.
\]
\[
d(3y)=3\,dy \;\Rightarrow\; 3\,dy\wedge dz.
\]
Summing: \(d\omega = (-2x+2x)\,dx\wedge dy + (-3+3)\,dy\wedge dz = 0\).

(b)
For a 1‑form \(\omega = \sum \omega_i dx^i\), the homotopy operator gives
\[
f(x) = \int_{0}^{1} \sum_i \omega_i(tx)\, x^i \, dt.
\]
Here \(\omega_x = 2xy\), \(\omega_y = x^{2}+3z\), \(\omega_z = 3y\). Substituting \(tx,ty,tz\):
\[
\omega_x(tx,ty,tz)=2t^{2}xy,\quad \omega_y = t^{2}x^{2}+3tz,\quad \omega_z = 3ty.
\]
Then
\[
\sum \omega_i x^i = 2t^{2}xy\cdot x + (t^{2}x^{2}+3tz)\cdot y + 3ty\cdot z = 3t^{2}x^{2}y + 6t y z.
\]
Integrate from \(t=0\) to \(1\):
\[
f = \int_{0}^{1} (3t^{2}x^{2}y + 6t y z)\,dt = x^{2}y + 3yz.
\]

(c)
\[
df = d(x^{2}y) + d(3yz) = (2xy\,dx + x^{2}\,dy) + (3z\,dy + 3y\,dz) = 2xy\,dx + (x^{2}+3z)\,dy + 3y\,dz = \omega.
\]
\end{solution}

\vfill\newpage

\begin{problem}
Let $U = \mathbb{R}^3$. Consider the 2-form $\omega = x \, dy \wedge dz + y \, dx \wedge dz$.
\begin{enumerate}
    \item[(a)] Compute $d\omega$ to verify that $\omega$ is closed.
    \item[(b)] Use the explicit homotopy operator formula $h\omega$ (for $k=2$) from the lecture:
    $$h\omega = \sum_{j=1}^k (-1)^{j-1} \left( \int_0^1 t^{k-1} f(tx) dt \right) x_{i_j} dx_{i_1} \dots \widehat{dx_{i_j}} \dots dx_{i_k}$$
    to find a 1-form $\eta = h\omega$ such that $d\eta = \omega$. \textit{(Hint: Apply the formula to $x \, dy \wedge dz$ and $y \, dx \wedge dz$ separately and add the results).}
    \item[(c)] Compute $d\eta$ directly from your answer in (b) to verify your primitive is correct.
\end{enumerate}
\end{problem}

\begin{solution}

(a)
\[
d\omega = dx\wedge dy\wedge dz + dy\wedge dx\wedge dz = dx\wedge dy\wedge dz - dx\wedge dy\wedge dz = 0.
\]

(b)
For a term \(f\,dx^{i}\wedge dx^{j}\), the formula is
\[
h(f\,dx^{i}\wedge dx^{j}) = \left(\int_{0}^{1} t f(tx)\,dt\right) \left(x^{i}dx^{j} - x^{j}dx^{i}\right).
\]
Apply to \(\omega_1 = x\,dy\wedge dz\): \(f=x\), \(i=2\), \(j=3\). Then
\[
\int_{0}^{1} t\cdot (tx)\,dt = \int_{0}^{1} t^{2}x\,dt = \frac{x}{3},
\]
so contribution: \(\frac{x}{3}(y\,dz - z\,dy) = \frac{xy}{3}dz - \frac{xz}{3}dy\).
Apply to \(\omega_2 = y\,dx\wedge dz\): \(f=y\), \(i=1\), \(j=3\). Then
\[
\int_{0}^{1} t\cdot (ty)\,dt = \frac{y}{3},
\]
contribution: \(\frac{y}{3}(x\,dz - z\,dx) = \frac{xy}{3}dz - \frac{yz}{3}dx\).
Summing:
\[
\eta = -\frac{yz}{3}dx - \frac{xz}{3}dy + \frac{2xy}{3}dz.
\]

(c)
\[
d\eta = d\left(-\frac{yz}{3}dx\right) + d\left(-\frac{xz}{3}dy\right) + d\left(\frac{2xy}{3}dz\right).
\]
\[
d\left(-\frac{yz}{3}dx\right) = -\frac{1}{3}(z\,dy + y\,dz)\wedge dx = \frac{z}{3}dx\wedge dy + \frac{y}{3}dx\wedge dz,
\]
\[
d\left(-\frac{xz}{3}dy\right) = -\frac{1}{3}(z\,dx + x\,dz)\wedge dy = -\frac{z}{3}dx\wedge dy + \frac{x}{3}dy\wedge dz,
\]
\[
d\left(\frac{2xy}{3}dz\right) = \frac{2}{3}(y\,dx + x\,dy)\wedge dz = \frac{2y}{3}dx\wedge dz + \frac{2x}{3}dy\wedge dz.
\]
Summing coefficients:
- \(dx\wedge dy\): \(\frac{z}{3} - \frac{z}{3} = 0\),
- \(dx\wedge dz\): \(\frac{y}{3} + \frac{2y}{3} = y\),
- \(dy\wedge dz\): \(\frac{x}{3} + \frac{2x}{3} = x\).
Thus \(d\eta = x\,dy\wedge dz + y\,dx\wedge dz = \omega\).
\end{solution}

\vfill

\end{document}